\section{Limitations and Future Work}
\label{app:limitations}

Section~\ref{sec:conclusion} lists the bounds on this study. Three of
them constrain the conclusions most, and each is taken up here with the
experiment that would lift it.

\paragraph{Single runs.}
Agent execution is stochastic, but each configuration was executed once. Large
gaps motivate the paper's conclusions. Small differences should be
treated as unresolved until replicated, among them the 0.2 to 1.7 pp between
\cond{Task-binary} and \cond{Checkpoint-binary} (Table~\ref{tab:exp2-full}) and the 0.8 to 1.6 pp between the best
fully-passed rates in Experiment~1 (Table~\ref{tab:exp1-full}). Stateful long-horizon execution makes exhaustive replication
costly, but cost does not remove the need to replicate the central contrasts.

\paragraph{Repeated-task distribution.}
R1 to R4 re-executed the same benchmark tasks. The experiment measures
deepening competence over recurring procedural workloads, not unrestricted
generalization to unseen distributions. Held-out task families and cross-domain
transfer are necessary to measure how broadly the learned procedures apply.

\paragraph{Human diagnosis.}
\cond{Checkpoint-binary + Human} provides the strongest feedback but requires human
effort after execution. Automating failure diagnosis without reverting to
unreliable self-judgment is the most direct route to scaling the approach.
